\section{Experimental Setup}
\label{sec:setup}

\paragraph{Benchmark and model.}
We use ALFWorld~\cite{shridhar2020alfworld}, the text household benchmark on which GACA reports its
ablations, with all six task types (pick-and-place, examine in light, clean, heat, cool, and pick two)
and a 50-step horizon. The policy is Qwen2.5-1.5B-Instruct~\cite{qwen2025qwen25}, the paper's smaller
scale, fully fine-tuned. Training runs in verl-agent~\cite{feng2025gigpo,sheng2025hybridflow}, the
stack GACA builds on, with the authors' ALFWorld prompt (ReAct-style~\cite{yao2022react} reasoning in
\texttt{<think>} tags, a command in \texttt{<action>} tags, 512-token responses, a two-step history) and
vLLM~\cite{kwon2023vllm} rollouts. Each method is a different advantage estimator plugged into the
same trainer; GRPO, GACA and all treatments call one shared, unit-tested implementation.

\paragraph{Recipe and budget.}
We keep GACA's shared configuration: group size 8, learning rate $10^{-6}$, $\gamma = 0.95$, KL
coefficient 0.01 with the low-variance estimator, invalid-format penalty 0.1, mean-and-std
normalization, $\lambda_0 = 0.5$, $\alpha = \beta = 0.5$. The budget is the one deliberate difference
(\cref{sec:budget}): on a single A100 40\,GB we train 4 tasks $\times$ 8 rollouts for a fixed 34
iterations, 1,088 trajectories per method, about 6\% of the published 19,200. All curves are still
rising at the end, so the comparison is of early learning at a matched rollout budget, and our absolute
success rates are not comparable to published ones.

\paragraph{Evaluation.}
Every method is evaluated every 10 iterations and at the end on the \emph{same} 128 held-out games
(ALFWorld's in-distribution validation split, one episode each, sampling temperature 0.4 as in the
authors' recipe). Because the games are shared, two methods are compared game by game with an exact
McNemar test~\cite{mcnemar1947} on the games one method solves and the other does not, which is far
more sensitive than comparing two rates at $n = 128$. We report the final checkpoint; the paper's
best-checkpoint convention picks the maximum of several noisy evaluations and is optimistic. One game
is 0.78 points. The base model solves 6 of 128 (4.7\%).

\paragraph{Methods.}
The controls are GRPO and GACA. The treatments are GACA with the $m_z$ reliability gate, GACA with the
leave-one-out score (GACA-LOO), COMMIT applied to GACA's step advantage, and two COMMIT variants
(\cref{sec:method}). One seed per method; the probe arm used a different training-environment seed
(the order in which training games are drawn) than the other arms, with the evaluation games pinned to
be identical. In an earlier pilot of this project, two seeds of the same method differed by 3.1 points at a matched
checkpoint, so differences of a few points between single runs are within seed noise.
